\documentclass[11pt, a4paper]{article}

\usepackage[a4paper, top=2cm, bottom=2cm, left=2.5cm, right=2.5cm]{geometry}
\usepackage[T1]{fontenc}
\usepackage[utf8]{inputenc}
\usepackage{microtype}
\usepackage{parskip}
\usepackage{hyperref}

\hypersetup{
    colorlinks=true,
    linkcolor=blue,
    filecolor=magenta,      
    urlcolor=blue,
}

\begin{document}

\begin{flushleft}
\textbf{Quoc Anh Trinh (Andrew)} \\
Melbourne, Victoria \\
+61 451 886 468 \\
\href{mailto:andrews.trinh@gmail.com}{andrews.trinh@gmail.com} \\
\href{https://www.linkedin.com/in/qatrinh/}{linkedin.com/in/qatrinh} \\
\vspace{1em}
\today
\end{flushleft}

\vspace{1em}

\begin{flushleft}
\textbf{To the AustralianSuper Hiring Team,}
\end{flushleft}

\vspace{1em}

As an AI Engineer specialising in agentic systems within the strictly regulated banking sector, I was immediately drawn to the \textbf{AI Engineer (Agentic)} role at AustralianSuper. I share your passion for building intelligence layers that go beyond task execution to enable true reasoning and planning. My background at Bendigo \& Adelaide Bank has been defined by exactly this: engineering the decision-making and reasoning layers that transform complex regulatory data into actionable insights for the Fund’s members.

In my current role within Financial Crime Risk, I am leveraging Agentic AI to bridge the gap between static regulation and active, high-performance monitoring. I recently engineered a \textbf{typology extractor} that utilized autonomous agents to transform 156 complex AUSTRAC guidance documents into 198 structured typologies. This project was not just about data extraction; it was an exercise in designing a reasoning framework that allows AI systems to interpret context, plan actions, and adapt to evolving regulatory landscapes.

My alignment with the Hyperautomation team’s mission is further evidenced by:
\begin{itemize}
    \item \textbf{Productionizing Intelligence:} I co-delivered an \textbf{AI-driven UAR-to-SMR pipeline} implementing secure \textbf{masking/rehydration} and \textbf{Pydantic-validated LLM prompting}; this reduced report prep time by \textbf{80\%} and shifted the focus toward strategic oversight.
    \item \textbf{RAG \& Tool Orchestration:} I have built production solutions using \textbf{Python and LlamaIndex}, designing domain-aware agent workflows for the \textbf{AUSTRAC Knowledge Hub}. This included architecting specialized agents for \textbf{typology extraction}, an \textbf{AML/CTF Program gap finder}, and a \textbf{conversational Q\&A interface} for querying complex regulatory logic with full auditability.
    \item \textbf{Safe-Autonomy in Regulated Environments:} Working in a "risk control-first" banking environment, I have developed a deep understanding of deploying AI with appropriate guardrails, governance, and auditability—ensuring \textbf{99.9\% logic alignment} and data integrity.
\end{itemize}

I am a builder who thrives in ambiguity and enjoys the abstract problem-solving required to design goal-driven decision logic. AustralianSuper’s commitment to delivering better outcomes for its 3 million members through innovation is a mission I am eager to support with my technical expertise in autonomous AI systems.

Thank you for your time and consideration. I look forward to the possibility of discussing how my experience in engineering reasoning agents for the financial sector can help AustralianSuper define the next generation of member-focused AI systems.

\vspace{2em}

Sincerely,

\vspace{1em}

Quoc Anh Trinh (Andrew)

\end{document}
